\documentclass[aspectratio=169]{beamer}

% =========================
% PAQUETES
% =========================
\usepackage[utf8]{inputenc}
\usepackage[T1]{fontenc}
\usepackage[spanish]{babel}
\usepackage{amsmath, amssymb}
\usepackage{graphicx}
\usepackage{booktabs}
\usepackage{tikz}
\usepackage{xcolor}

% =========================
% TEMA
% =========================
\usetheme{Madrid}
\usecolortheme{default}
\setbeamertemplate{navigation symbols}{}
\setbeamertemplate{footline}{%
    \hfill
    \raisebox{1.4ex}[0pt][0pt]{%
        \scriptsize\insertframenumber\,/\,\inserttotalframenumber\hspace*{1ex}%
    }%
    \vspace{1.5ex}%
}

% =========================
% INFORMACIÓN
% =========================
\title{Localización de una Fuente Sísmica \\
mediante un Problema Inverso y Datos Simulados}

\subtitle{Proyecto de Cálculo Multivariado}

\author{
Kevin Buesaquillo \\
Danilo Lara \\
David Cabrera \\
Alejandro Parra \\
Luis Diaz
}

\institute{
Ingeniería de Software
}

\date{2026}

% =========================
% DOCUMENTO
% =========================
\begin{document}

% =========================
% PORTADA
% =========================

\begin{frame}
    \titlepage
\end{frame}

% =========================
% DIAPO DANILO
% =========================

\begin{frame}{Danilo — Introducción y Modelo Matemático}

\textbf{Rol en la exposición:}

Introducir el problema sísmico y explicar el modelo matemático base.

\vspace{0.4cm}

\textbf{Todo lo que debe decir:}

\begin{itemize}
    \item Explicar qué se busca en el proyecto.
    \item Definir qué es una fuente sísmica.
    \item Explicar qué representan los sensores.
    \item Explicar la idea general del problema inverso.
    \item Diferenciar problema directo vs problema inverso.
    \item Explicar por qué se usa cálculo multivariado.
    \item Introducir el modelo de atenuación sísmica.
    \item Explicar qué representa la amplitud inicial $A_0$.
    \item Explicar qué representa la distancia $R_i$.
    \item Explicar por qué la amplitud disminuye con la distancia.
    \item Introducir la fórmula:
    
    $$
    A_{pred,i}=A_0\frac{e^{-R_i}}{R_i}
    $$

    \item Explicar que el objetivo es encontrar la ubicación de la fuente.
    \item Conectar el proyecto con funciones de varias variables.
\end{itemize}

\end{frame}

% =========================
% DIAPO LUIS
% =========================

\begin{frame}{Luis — Simulación de Datos y Sensores}

\textbf{Rol en la exposición:}

Explicar cómo se generaron los datos simulados utilizados en el proyecto.

\vspace{0.4cm}

\textbf{Todo lo que debe decir:}

\begin{itemize}
    \item Explicar que esta parte corresponde al problema directo.
    \item Presentar la fuente sísmica real simulada.
    \item Explicar las coordenadas de la fuente.
    \item Explicar por qué $z$ es negativo.
    \item Explicar el valor usado para $A_0$.
    \item Explicar cómo se distribuyeron los sensores.
    \item Explicar por qué se utilizaron 12 sensores.
    \item Mostrar la distribución espacial de los sensores.
    \item Explicar la distancia euclidiana en $\mathbb{R}^3$.
    \item Introducir la fórmula:
    
    $$
    R_i = \sqrt{(x_i-x)^2+(y_i-y)^2+(z_i-z)^2}
    $$

    \item Explicar cómo se calcularon las amplitudes ideales.
    \item Explicar qué es el ruido gaussiano.
    \item Explicar por qué se agregó ruido.
    \item Explicar cómo se obtuvieron las amplitudes observadas.
    \item Explicar que los sensores cercanos registran amplitudes mayores.
    \item Explicar que estos datos serán la entrada del problema inverso.
\end{itemize}

\end{frame}

% =========================
% DIAPO DAVID
% =========================

\begin{frame}{David — Función de Error y Optimización}

\textbf{Rol en la exposición:}

Explicar cómo se construyó la función de error y cómo se encontró el mínimo.

\vspace{0.4cm}

\textbf{Todo lo que debe decir:}

\begin{itemize}
    \item Explicar qué significa comparar amplitudes observadas y predichas.
    \item Explicar por qué se construye una función de error.
    \item Explicar qué significa minimizar el error.
    \item Explicar por qué se usan mínimos cuadrados.
    \item Introducir la función:
    
    $$
    E(x,y,z)=\sum_{i=1}^{n}(A_{obs,i}-A_{pred,i})^2
    $$

    \item Explicar que el mínimo representa la mejor estimación.
    \item Explicar el método de optimización Nelder-Mead.
    \item Explicar que también se utilizó búsqueda por grilla como apoyo visual.
    \item Explicar qué ocurre cuando un punto candidato está cerca de la fuente real.
    \item Presentar la fuente estimada.
    \item Comparar fuente real vs fuente estimada.
    \item Explicar el significado físico del error pequeño obtenido.
    \item Explicar qué son los residuales.
    \item Explicar por qué residuales pequeños indican buen ajuste.
\end{itemize}

\end{frame}

% =========================
% DIAPO ALEJANDRO
% =========================

\begin{frame}{Alejandro — Visualización y Mapas de Calor}

\textbf{Rol en la exposición:}

Explicar las visualizaciones de la función de error y la interpretación gráfica.

\vspace{0.4cm}

\textbf{Todo lo que debe decir:}

\begin{itemize}
    \item Explicar por qué la función $E(x,y,z)$ no se puede visualizar completa fácilmente.
    \item Explicar qué significa fijar un valor de $z$.
    \item Explicar los cortes bidimensionales.
    \item Explicar qué representan los mapas de calor.
    \item Explicar qué representan las curvas de nivel.
    \item Explicar cómo identificar regiones de bajo error.
    \item Explicar el corte en $z=-1.10$.
    \item Explicar el corte en $z=-1.14$.
    \item Explicar por qué esos cortes son importantes.
    \item Explicar cómo cambia el error cuando cambia la profundidad.
    \item Explicar la gráfica $E_{min}(z)$.
    \item Explicar la trayectoria del mínimo en el plano $(x,y)$.
    \item Explicar la animación de cortes en profundidad.
    \item Explicar que la visualización confirma el resultado numérico.
\end{itemize}

\end{frame}

% =========================
% DIAPO KEVIN
% =========================

\begin{frame}{Kevin — Conclusiones y Defensa Final}

\textbf{Rol en la exposición:}

Cerrar el proyecto, interpretar los resultados y responder preguntas finales.

\vspace{0.4cm}

\textbf{Todo lo que debe decir:}

\begin{itemize}
    \item Explicar el resultado final del proyecto.
    \item Explicar que la fuente estimada quedó cerca de la fuente real.
    \item Explicar qué significa físicamente el mínimo encontrado.
    \item Explicar por qué el modelo es coherente.
    \item Explicar la relación entre distancia y amplitud.
    \item Explicar cómo interviene el cálculo multivariado.
    \item Explicar la importancia de las funciones de varias variables.
    \item Explicar la importancia de la optimización.
    \item Explicar las limitaciones del modelo.
    \item Explicar posibles mejoras futuras.
    \item Concluir que el problema inverso pudo resolverse correctamente.
    \item Cerrar la exposición de manera formal.
\end{itemize}

\end{frame}

% =========================
% Limitaciones
% =========================

\begin{frame}{Limitaciones y Mejoras Futuras}

\textbf{Rol en la exposición:}

Explicar las limitaciones del modelo y posibles mejoras para futuros trabajos.

\vspace{0.4cm}

\textbf{Todo lo que debe decir:}

\begin{itemize}
    \item Explicar que el modelo de atenuación es simplificado.
    \item Explicar que no se consideraron efectos de propagación complejos.
    \item Explicar que el ruido real puede no ser perfectamente gaussiano.
    \item Explicar que se asumió un medio homogéneo.
\end{itemize}

\end{frame}


% =========================
% DIAPO FINAL
% =========================

\begin{frame}{Gracias}

\centering

{\Huge Gracias por su atención}

\vspace{1cm}

Preguntas

\end{frame}

\end{document}
```
